\section*{Discussion}

The preregistered transport-failure hypothesis was not confirmed once a
verified estimator defect was corrected. The original analysis counted accepted
studies whose report mentioned none of the five targets as correctly
classified. These studies made up about half of the source cohort and a fifth
of the external one, so the originally reported transfer gaps measured a difference in
label availability between two report corpora rather than a change in model
error. The near-zero corrected gaps are themselves fragile: label handling,
label source, aggregation, view and training seed each moved them by more than
their intervals, in both directions. The original question, whether a frozen
selective policy's error rises at a new hospital, cannot be answered with these
labels.

Three findings were more stable. First, frozen thresholds kept coverage but not
operating points: external specificity fell for several findings under both
label handlings examined. A deployment monitoring coverage alone would not
detect such a shift. Second, misaligned context harmed late fusion more than
absent context at the source under every specification, permutation and seed.
The same effect at the external site did not survive the second seed. Third,
better error ranking at the source, from a decision-aware score, did not make
the operating point more transportable.

\subsection*{Is the correction outcome-driven?}
The correction was made after the external results were known, and it reverses
the preregistered primary result, so it deserves scrutiny. Four features argue
against an outcome-driven choice. First, the defect contradicts the
implementation's own documented intent, and the corrected estimator applies the
registered definition without other changes. Second, the original estimator is
an exact multiple of the corrected one, and the multiplier, the evaluable share
of accepted studies, is reported for both sites. Third, the correction removes
a confirmed hypothesis rather than creating one. Fourth, no alternative
aggregation or label handling restores a robust transfer failure: the sign
depends on the choice. We report both estimators and every alternative so that
readers can weigh this history.

\subsection*{What the endpoint measures}
Selective error over report-derived labels is an error rate over the cells a
labeller chose to label, in the studies a policy accepted. Three mechanisms can
change it between hospitals. Case mix differed, with effusion and edema roughly
twice as prevalent externally when unmentioned findings counted as negative.
Documentation differed: source impressions often omit all targets, external
impressions carry more uncertain mentions, and findings sections are absent
from most external studies. Labeller accuracy may also differ, because report
labels are imperfect image proxies and the external labelling settings are
undocumented.\cite{jain2021} Matching aggregate positive rates therefore did
not establish comparable labels. The controlled analysis shows that selective
missingness, meaning which cells carry a label at all, drove the originally reported
contrast and remains a first-order driver after correction. This sharpens an
earlier argument that cross-dataset failure of chest-radiograph models reflects
label shift.\cite{cohen2020} Here, label shift is not a competing explanation
for model behaviour but a property of the measurement itself. It is strong
enough to set the sign of a selective-risk comparison while predictions and
thresholds are held fixed.

\subsection*{Implications for evaluation practice}
The defect is easy to make and invisible in aggregate output. Any per-study
loss averaged over observed labels needs an explicit rule for studies with
none; here a guard against division by zero silently chose one. For cross-site
evaluations built on report-derived labels, we recommend six practices:
\begin{itemize}
  \item report, per site, the share of studies and of accepted studies with at
        least one labelled target;
  \item test how the estimator treats unlabelled studies;
  \item fix and report the handling of uncertain and unmentioned labels;
  \item report pathology-specific operating characteristics;
  \item replay context interventions to confirm their unit and constraints;
  \item quantify calibration-sample uncertainty alongside the evaluation
        bootstrap.
\end{itemize}
Where the sign of a cross-site contrast matters, a reference standard that is
dense and comparable at both sites is needed. We did not evaluate target-site
recalibration, which the protocol prohibits. It answers a different question,
whether a method can be re-tuned, and we make no claim about its effect.

\subsection*{Limitations}
All labels are automatic and report-derived. No expert reference standard was
available; a blinded adjudication protocol is provided (\ssec{ssec:adjud}). The
estimator correction and all revision analyses followed inspection of external
results. The correction fixes a verified defect but reverses the original conclusion, and both estimates are reported. Intervals condition on one
calibration sample and one C2 realisation. Two training seeds cannot establish
seed-independence, and the second seed changed several estimates materially.
C2 was applied per image; synthetic context removal differs from naturally
missing context; and the informativeness rule measures length, not content.
The evaluation covers one pair of hospitals in one direction. Demographic
analyses are external-only and exploratory. Two pre-data planning documents
existed; the prose plan's one-image-per-study and concordant-study
sensitivities were not run.
